% Generated draft: requires mathematical and source review.
\documentclass[11pt]{article}
\usepackage[margin=25mm]{geometry}
\usepackage{fontspec}
\setmainfont{Latin Modern Roman}
\usepackage{amsmath,amssymb,mathtools,unicode-math}
\setmathfont{Latin Modern Math}
\usepackage{enumitem,needspace,xurl,hyperref}
\setlength{\parindent}{0pt}
\setlength{\parskip}{5pt}
\setlength{\emergencystretch}{3em}
\newcommand{\N}{\mathbb{N}}
\newcommand{\Z}{\mathbb{Z}}
\newcommand{\Q}{\mathbb{Q}}
\newcommand{\R}{\mathbb{R}}
\newcommand{\C}{\mathbb{C}}
\newcommand{\F}{\mathbb{F}}
\newcommand{\A}{\mathbb{A}}
\newcommand{\PP}{\mathbb{P}}
\newcommand{\E}{\mathbb{E}}
\newcommand{\eps}{\varepsilon}
\newcommand{\dd}{\,\mathrm{d}}
\begin{document}
\section*{UnsolvedMath ID: 30005738}
\textbf{Source problem code:} catalogue entry
\subsection{Definitions and mathematical statement}
Let \(k\) be a field. The supplied record does not specify its characteristic, whether it is algebraically closed, or any unit assumptions. These restrictions must be recovered before a definitive formulation can be given.

Under the standard algebraic convention, a transposed Poisson algebra (TPA) is a vector space \(A\) over \(k\), equipped with a commutative associative bilinear multiplication \(\cdot\) and a Lie bracket \([\ ,\ ]\), satisfying
\[
2z\cdot[x,y]=[z\cdot x,y]+[x,z\cdot y]
\qquad(x,y,z\in A).
\]
Here a Lie bracket is alternating and satisfies the Jacobi identity. This definition supplies terminology; the record does not establish the intended conventions in exceptional characteristics.

For the standard definition of specialness stated here, suppose \(\operatorname{char}k\ne2\). A Jordan algebra is a commutative algebra \((J,\circ)\) satisfying
\[
(x^2\circ y)\circ x=x^2\circ(y\circ x),
\qquad x^2=x\circ x,
\qquad x,y\in J.
\]
It is special if there exist an associative \(k\)-algebra \(B\) and an injective \(k\)-linear map \(\iota:J\to B\) such that
\[
\iota(x\circ y)=\frac{\iota(x)\iota(y)+\iota(y)\iota(x)}{2}
\qquad(x,y\in J).
\]
No unit or finite-dimensionality condition on \(B\) is imposed by this standard definition. The characteristic restriction here is an explicit convention for this definition, not a hypothesis recovered from the source.

The supplied reconstruction asks to characterize the Jordan brackets arising from TPAs and to determine whether their associated Jordan algebras are special. The quoted original asks, ``Are these Jordan algebras special?'' However, neither the meaning of Jordan bracket, the construction of the associated Jordan algebra, nor the antecedent of ``these'' is supplied. Thus the first question can only be recorded conditionally: for every Jordan algebra in the intended TPA-derived family, does an embedding \(\iota\) as above exist? The family itself and the requested characterization remain unspecified; no replacement construction is asserted here.

The second question asks whether there exists a finite-dimensional simple TPA whose two multiplications are both simple. To display the relevant ideal conditions, for a subspace \(I\subseteq A\) write
\[
\begin{aligned}
I\text{ is an associative ideal}&\iff A\cdot I\subseteq I,\\
I\text{ is a Lie ideal}&\iff [A,I]\subseteq I,\\
I\text{ is a common ideal}&\iff A\cdot I\subseteq I\ \text{and}\ [A,I]\subseteq I.
\end{aligned}
\]
Under the interpretation that TPA simplicity concerns common ideals and that simplicity of each multiplication concerns its own ideals, the requested existence conditions include
\[
0<\dim_k A<\infty,
\]
\[
\forall I\le_k A,\quad
A\cdot I\subseteq I\ \Longrightarrow\ I\in\{0,A\},
\qquad
\forall I\le_k A,\quad
[A,I]\subseteq I\ \Longrightarrow\ I\in\{0,A\}.
\]
These conditions also exclude proper nonzero common ideals. The source does not specify whether simplicity additionally requires \(A\cdot A\ne0\), \([A,A]\ne0\), or a dimension restriction excluding one-dimensional cases. Accordingly, the displayed conditions record the ideal requirements but do not settle those missing conventions or the intended ground field. The question remains a statement-only draft requiring mathematical review.
\subsection{Short English statement}
The record asks which Jordan brackets come from transposed Poisson algebras, whether the associated Jordan algebras are special, and whether a finite-dimensional simple transposed Poisson algebra can have both operations simple. The Jordan construction, ground field, and precise simplicity conventions are missing, so a fully determined statement cannot yet be extracted.
\subsection{Sources}
\begin{itemize}
\item Supplied local UnsolvedMath registry record, permanent ID \(30005738\), title: \emph{Jordan Brackets and Simple Transposed Poisson Algebras}; exact source problem code: \texttt{catalogue entry}. The record belongs to the Oberwolfach Reports Open Problems collection and supplies no named proposer.
\item Supplied authoritative source citation: \emph{Mini-Workshop: Poisson and Poisson-type algebras} (2024), Oberwolfach Reports, pp.\ 2681--2715, DOI \texttt{10.4171/owr/2023/46}: \url{https://doi.org/10.4171/owr/2023/46}. The original excerpt identifies p.\ 2697. The supplied literature-source label is \emph{Oberwolfach Report 46/2023}, with the same link. The record also gives proposed year 2023; the differing year designations are retained without reconciliation.
\item The supplied record labels the statement \texttt{corrected\_verified} and says that it was reconstructed and checked against the official report. This is supplied verification metadata; no independent source inspection is claimed here.
\item Literature assessment dated exactly \texttt{2026-08-22}: the supplied curation labels the problem open and states that known results rule out complex simple finite-dimensional TPAs in relevant cases and provide formal-series examples, while the finite-dimensional example with both multiplications simple remains posed as open. The relevant cases, hypotheses, and example constructions are not specified. This is supplied dated curation, not an independently verified literature assessment; its status and cited claims require independent checking.
\end{itemize}
\end{document}
